% ==============================================================================
% Plantilla Institucional de Trabajo Terminal — UPIIZ IPN
% Archivo: ejemplos/referencias.tex
% ==============================================================================
% Guía práctica para citación bibliográfica en formato IEEE.
% ==============================================================================

\section{Guía de ejemplos para citación bibliográfica IEEE}
\label{sec:ejemplo-citas}

El formato IEEE utiliza citas numéricas correlativas entre corchetes rectos \texttt{[1]}, \texttt{[2]}, generadas automáticamente mediante el comando \texttt{\textbackslash cite\{clave\}}.

A continuación se muestran ejemplos de citación para los distintos tipos de fuentes bibliográficas:

\begin{itemize}
    \item \textbf{Libro de texto:} Como se describe en los fundamentos de ingeniería de control moderno por Ogata \cite{ogata2010modern}.
    \item \textbf{Artículo de revista indexada (Journal):} Conforme a las metodologías de control difuso documentadas por Smith et al. \cite{smith2022control}.
    \item \textbf{Artículo de conferencia o congreso:} De acuerdo con los algoritmos de estimación de estados presentados en ICRA \cite{ieee2021robotics}.
    \item \textbf{Hoja técnica o manual de fabricante (Datasheet):} Según las especificaciones de consumo y periféricos del microcontrolador \cite{arizona2023embedded}.
    \item \textbf{Norma internacional o estándar técnico:} Cumpliendo con los lineamientos de seguridad en sistemas de control de maquinaria \cite{iso2018safety}.
    \item \textbf{Documento institucional o sitio web:} Atendiendo a la normativa y reglamentación del Instituto Politécnico Nacional \cite{ipn2025normativa}.
    \item \textbf{Citas múltiples:} Se pueden agrupar varias fuentes en un solo comando \texttt{\textbackslash cite\{ogata2010modern, smith2022control, ieee2021robotics\}}, produciendo la referencia compacta \cite{ogata2010modern, smith2022control, ieee2021robotics}.
\end{itemize}
